\documentclass[12pt]{article}
\usepackage{amsmath,amssymb,amsfonts}
\usepackage[margin=1in]{geometry}

\begin{document}

\section*{Chapter 2, Problem 6}

\subsection*{Problem Statement}
Show that if $f$ also involves the second derivative of $y$ with respect to $x$, i.e., $f = f(x, y, y', y'')$, then for fixed endpoints and prescribed $y'$ at the endpoints, the Euler--Lagrange equation is
\[
  \frac{\partial f}{\partial y} - \frac{d}{dx}\frac{\partial f}{\partial y'} + \frac{d^2}{dx^2}\frac{\partial f}{\partial y''} = 0.
\]
What continuity assumptions must be made about derivatives of $y$?

\subsection*{Solution}

Consider the functional
\[
  I[y] = \int_{x_1}^{x_2} f(x, y, y', y'')\,dx.
\]
Let $y(x)$ be the extremizing function and consider a variation $y(x) + \epsilon\,\eta(x)$, where $\eta(x)$ is an arbitrary function satisfying the boundary conditions:
\[
  \eta(x_1) = \eta(x_2) = 0, \qquad \eta'(x_1) = \eta'(x_2) = 0.
\]
The first condition fixes the endpoints; the second fixes the prescribed values of $y'$ at the endpoints.

The first variation is:
\[
  \frac{dI}{d\epsilon}\bigg|_{\epsilon=0}
  = \int_{x_1}^{x_2} \left[\frac{\partial f}{\partial y}\,\eta + \frac{\partial f}{\partial y'}\,\eta' + \frac{\partial f}{\partial y''}\,\eta''\right] dx = 0.
\]

We integrate the second term by parts once:
\[
  \int_{x_1}^{x_2} \frac{\partial f}{\partial y'}\,\eta'\,dx
  = \left[\frac{\partial f}{\partial y'}\,\eta\right]_{x_1}^{x_2}
  - \int_{x_1}^{x_2} \frac{d}{dx}\frac{\partial f}{\partial y'}\,\eta\,dx.
\]
The boundary term vanishes since $\eta(x_1) = \eta(x_2) = 0$.

We integrate the third term by parts twice:
\[
  \int_{x_1}^{x_2} \frac{\partial f}{\partial y''}\,\eta''\,dx
  = \left[\frac{\partial f}{\partial y''}\,\eta'\right]_{x_1}^{x_2}
  - \int_{x_1}^{x_2} \frac{d}{dx}\frac{\partial f}{\partial y''}\,\eta'\,dx.
\]
The boundary term vanishes since $\eta'(x_1) = \eta'(x_2) = 0$. Integrating by parts again:
\[
  -\int_{x_1}^{x_2} \frac{d}{dx}\frac{\partial f}{\partial y''}\,\eta'\,dx
  = -\left[\frac{d}{dx}\frac{\partial f}{\partial y''}\,\eta\right]_{x_1}^{x_2}
  + \int_{x_1}^{x_2} \frac{d^2}{dx^2}\frac{\partial f}{\partial y''}\,\eta\,dx.
\]
The boundary term again vanishes since $\eta(x_1) = \eta(x_2) = 0$.

Collecting all terms:
\[
  \int_{x_1}^{x_2} \left[\frac{\partial f}{\partial y} - \frac{d}{dx}\frac{\partial f}{\partial y'} + \frac{d^2}{dx^2}\frac{\partial f}{\partial y''}\right] \eta\,dx = 0.
\]

Since $\eta(x)$ is arbitrary (subject to the boundary conditions), the fundamental lemma of the calculus of variations gives:
\[
  \boxed{\frac{\partial f}{\partial y} - \frac{d}{dx}\frac{\partial f}{\partial y'} + \frac{d^2}{dx^2}\frac{\partial f}{\partial y''} = 0.}
\]

\medskip
\noindent\textbf{Continuity assumptions:} The extremizing function $y(x)$ must have continuous derivatives up to fourth order ($y \in C^4$). This is because:
\begin{itemize}
  \item $\partial f/\partial y''$ generally involves $y''$, so its second derivative $\frac{d^2}{dx^2}\frac{\partial f}{\partial y''}$ involves $y^{(4)}$.
  \item The Euler--Lagrange equation must hold pointwise, which requires all the derivatives appearing in it to exist and be continuous.
\end{itemize}

\end{document}
